\section{Prompts}
\label{app:prompts}
\ifconf
Fig.~\ref{fig:minprompt} gives the minimal prompt used by the \emph{no method} ablation. The investigation-method
prompt of the full configuration (81 lines) is in the artifact (\texttt{agent/robust/prompt.py}); its parts are
summarised in \S\ref{sec:method}, and the instruction discussed in \S\ref{sec:study} reads ``Do not look up
what you already know: the host and user named in the alert, the relation that only says a process runs on
that host''. Field names refer to the aliases presented to the model (\S\ref{sec:alias}).
\else
Fig.~\ref{fig:prompt} reproduces the investigation-method system prompt used by the full configuration, and
Fig.~\ref{fig:minprompt} the minimal prompt used by the \emph{no method} ablation. Both are taken verbatim from
the released code (\texttt{agent/robust/prompt.py}). Field names refer to the aliases presented to the model
(\S\ref{sec:alias}).

\begin{figure*}[!t]
\centering
\begin{minipage}{0.92\textwidth}
\VerbatimInput[fontsize=\scriptsize, baselinestretch=\PromptStretch, frame=single, framesep=2mm]{prompts/method_prompt.txt}
\end{minipage}
\caption{Investigation-method system prompt (full configuration).}
\label{fig:prompt}
\end{figure*}
\fi

\begin{figure*}[!t]
\centering
\begin{minipage}{0.92\textwidth}
\VerbatimInput[fontsize=\scriptsize, frame=single, framesep=2mm]{prompts/minimal_prompt.txt}
\end{minipage}
\caption{Minimal system prompt (\emph{no method} ablation).}
\label{fig:minprompt}
\end{figure*}

\ifconf
\section{Artifact}
\label{app:repro}
The anonymised artifact, \ArtifactURL, contains the agent, the benchmark and its variants, the traces of
every run, and the scripts that regenerate every table, figure and number in this paper byte for byte; its
README gives the commands. The agent's tests run against the competition's public reference client.
\else
\section{Reproducing the tables}
\label{app:repro}
Every table and plot in \S\ref{sec:benchmark} and \S\ref{sec:results} is generated from the released files:
\begin{itemize}
\item \texttt{scripts/benchmark\_to\_tex.py} writes the task and type tables and the benchmark statistics from
the benchmark zip;
\item \texttt{scripts/results\_to\_tex.py} writes the pilot, ablation, breakdown and per-task tables, the
figure data and the numbers quoted in the text from the run folders.
\end{itemize}
\fi

\section{Engineering notes}
\label{app:engineering}
\emph{Policy-compatible rationale fields.} The runtime requires a free-text field named \texttt{reasoning} on
every evidence call and in the submission. In early runs with a recent frontier model, the first request of
every task was declined before any output, with the refusal category \texttt{reasoning\_extraction}: a prompt
asking the model to ``state its hypotheses'' together with a required field called \texttt{reasoning} described
as the ``investigation trace'' read as a request to externalise the model's internal reasoning, which the
provider's usage policy declines. The fix keeps the runtime contract intact: the instruction to state
hypotheses was removed from the first message (the method still asks the model to consider them), the agent
presents the field to the model as \texttt{purpose} (``one short sentence naming what this call checks'') and
the submission's field as \texttt{summary} (``a short incident summary''), and maps both back to
\texttt{reasoning} before validation and sending; the model's actual reasoning is read from its summarised
thinking blocks. After the change, all tasks of a mock run completed where all had previously been refused. Any
agent built for a fixed runtime contract can meet the same conflict, and a field rename is enough to resolve it.

\emph{Parallel execution with a learned limit.} Budget-scored evaluations are often time-boxed. The agent runs
$W$ workers that each lease and investigate tasks. Before leasing, a worker reserves a slot, so no more leases
are attempted than the current limit allows. If the runtime refuses a lease because of its active-task limit
while other workers are busy, the pool lowers its limit to the number of tasks in flight and the worker waits
for a free slot; refusals while nothing else runs are waited out with a bounded number of retries. Short
operation-capacity or rate-limit refusals on tool calls are retried with backoff; a refused call was not
executed, so retrying spends nothing. One task's failure is contained to that task.

\ifconf
% Ethics considerations and generative AI usage (IEEE S&P / SaTML requirements).
% In the body of the journal build, in the appendix of the conference build.
\section{Ethics Considerations}
\label{sec:ethics}
\paragraph{Data and human subjects} \bench is derived solely from public, permissively licensed recordings made
by their publishers in lab environments (OTRF Security-Datasets, MIT; Splunk attack\_data, Apache-2.0). They
contain lab host and account names and no personal data; no human subjects were involved and no new data was
collected.

\paragraph{Live systems} No live system was measured, tested or accessed. All investigations ran against a
local runtime serving the recorded graphs; the only external service used was the model provider's API, under
its terms of use.

\paragraph{Tools and attacks} The agent is defensive: it reads telemetry and recommends responses, and takes no
action on systems. The mimicry variant shows that renaming attacker-chosen artefacts can lead an LLM
investigation agent to clear real attacks. Renaming tools and accounts is long-standing intruder tradecraft,
so publishing the result gives attackers little they lack, while defenders learn not to trust names and how to
measure that. Mitigations follow from the analysis: require behaviour-based justification for clearing an
alert, check grounding against identifier types, and evaluate on cue-controlled data. The finding concerns how
investigation agents weigh names, not a flaw in a specific deployed product, so no vulnerability disclosure
process applied.

\paragraph{Availability} Code (MIT), the benchmark and its variants (with the upstream licences), the traces of
every tool call, and the scripts that regenerate every table are released with the paper\ifconf; an
anonymised version for review is at \ArtifactURL\fi.

\section{Generative AI Usage}
\label{sec:genai}
\paragraph{As object of study} The agent under study is driven by Claude Opus 5.5 (\texttt{claude-opus-5-5})
through the Anthropic API, and, for the model comparison, by GPT-6 astra (\texttt{gpt-6-astra}) through the
OpenAI Responses API. The models are closed, and later versions may behave differently; we pin the model
identifiers, run three repeats per condition to measure variance, and release the trace of every call so that
results can be inspected without re-running the model.

\paragraph{Footprint} Across all recorded runs -- the ablation study, the cue-controlled, three-model and RQ7
remedy experiments, and the preliminary runs that preceded them -- the agent ran 82 times over 1{,}677 task
investigations, making 10{,}784 model calls with 2.6\,M output tokens, 51.0\,M prompt tokens served from the
providers' caches and 11.1\,M uncached input tokens, across three models (Claude Opus 5.5, GPT-6 astra and
DeepSeek-V3.2; a weaker open model, Llama~3.3~70B, was piloted and discarded). The per-call accounting of every
run is in the released traces. Queries were kept low by prompt caching, a per-task call budget, a small
benchmark, and offline rescoring: when the ground truth was revised, earlier runs were rescored from their traces
instead of being repeated. Experiments ran on the providers' hosted inference; local computation (graph building,
scoring, the agent loop) ran on a small cloud server.

\paragraph{In preparing this work} % AUTHOR: confirm this matches your actual division of labour before submitting.
The authors conceived and directed the study, designed the benchmark, specified the ground-truth rule, ran all
model experiments and verified the results. A general-purpose LLM-based coding assistant was used as a tool
throughout: to write and refactor the agent, the benchmark and its variant builders and the analysis scripts;
to set up, run and debug the experiment harnesses; and to draft and edit this manuscript. Every reported number
is produced from the released traces by the released scripts; the authors validated all generated code, text and
figures and are solely responsible for the content. The assistant generated no experimental result or
measurement.


\section{Supporting tables}
\label{app:tables}
Table~\ref{tab:tasks} lists the tasks of \bench. The tool-interface counts for the study and the scoring
sensitivity analysis are generated by \texttt{scripts/results\_to\_tex.py} and included in the artifact to
save space here.
\begin{table*}[t]\centering\scriptsize
\caption{The 21 tasks of \bench. Graph size is entities/relations; \emph{Core} is the number of core entities (malicious) or legitimacy anchors (benign) in the ground truth.}
\label{tab:tasks}
\begin{tabular}{@{}lllllrr@{}}\toprule
Task & OS & Verdict & ATT\&CK & Alert title & Graph & Core\\\midrule
\texttt{lsass-dump-comsvcs} & Win & malicious & T1003.001 & Process opened LSASS with full access & 88/115 & 5\\
\texttt{logon-script-persistence} & Win & malicious & T1037.001 & User environment registry value written by reg.exe & 47/75 & 6\\
\texttt{service-binpath-hijack} & Win & malicious & T1543.003 & Service configuration changed with sc.exe & 157/219 & 6\\
\texttt{hta-startup-folder} & Win & malicious & T1218.005 & mshta.exe executed an HTA file & 338/717 & 6\\
\texttt{firewall-rule-added} & Win & malicious & T1562.004 & Inbound firewall rule added from the command line & 174/365 & 4\\
\texttt{python-http-server} & Win & malicious & T1059 & Interpreter from a user profile started a network service & 300/633 & 6\\
\texttt{bits-download} & Win & malicious & T1197 & BITS transfer job started from the command line & 186/247 & 4\\
\texttt{run-key-and-beacon} & Win & malicious & T1547.001 & Autostart Run key value set by PowerShell & 74/87 & 5\\
\texttt{lsass-query-by-svchost} & Win & benign & - & Process opened a handle to LSASS & 186/247 & 2\\
\texttt{w32time-runtime-key} & Win & benign & - & Registry value written under a path matching '\textbackslash{}Run' & 300/633 & 3\\
\texttt{gpsvc-service-host} & Win & benign & - & New service host process started as SYSTEM & 174/365 & 2\\
\texttt{linux-shadow-read} & Lin & malicious & T1003.008 & Credential file read by a script run with sudo & 50/76 & 6\\
\texttt{linux-wiper-shred} & Lin & malicious & T1485 & shred overwrote a system directory & 317/672 & 7\\
\texttt{linux-kernel-module} & Lin & malicious & T1547.006 & Kernel module loaded with insmod & 78/142 & 6\\
\texttt{linux-account-created} & Lin & malicious & T1136.001 & Local account created & 76/127 & 5\\
\texttt{linux-sudoers-nopasswd} & Lin & malicious & T1548.003 & Passwordless sudo rule written & 141/245 & 5\\
\texttt{linux-ld-preload-hijack} & Lin & malicious & T1574.006 & Freshly compiled binary executed as root from a home directory & 150/265 & 8\\
\texttt{linux-service-stopped} & Lin & malicious & T1489 & Service state changed by systemctl as root & 1602/3951 & 5\\
\texttt{linux-dpkg-service-start} & Lin & benign & - & Service state changed by systemctl as root & 1602/3951 & 3\\
\texttt{linux-motd-discovery} & Lin & benign & - & Session discovery command run as root & 317/672 & 3\\
\texttt{linux-motd-tmp-cleanup} & Lin & benign & - & Files deleted by rm as root & 317/672 & 3\\
\bottomrule\end{tabular}\end{table*}

\fi
